%! Author = abenjeddou
%! Date = 30/06/2026

\section{Introduction}
After completing the study, design, and architecture of our system, we move on to the implementation sections.
This chapter marks the beginning of the development phase, organized into \textbf{6 sprints} following the Agile Scrum methodology with four-week iterations.
Each sprint delivers a functional increment of the system, from the initial bare project setup through to full observability in production.

Table 4.1 provides an overview of the sprint planning.

\begin{center}
    \justifying
    \begin{tabularx}{\textwidth}{|c | p{2cm} | p{3cm} | X|}
\toprule
\textbf{Sprint} & \textbf{Period} & \textbf{Theme} & \textbf{Key Deliverables} \\
\midrule
1 & Dec 01 -- Dec 28 & Project Setup \& Data Source & Maven module structure, development environment, RabbitMQ source operator \\
\hline
2 & Dec 29 -- Jan 25 & Message Validation \& Subscription & Message validation, Valkey integration, main record retrieval, rule validation \\
\hline
3 & Jan 26 -- Feb 22 & Notification Delivery & Erable API call, retry mechanism, state update sink \\
\hline
4 & Feb 23 -- Mar 22 & AI Prediction Module & Feature extraction, ONNX inference, rule-based fallback \\
\hline
5 & Mar 23 -- Apr 19 & Containerization \& Deployment & Docker image, Kubernetes, CI/CD pipeline \\
\hline
6 & Apr 20 -- May 17 & Observability \& Monitoring & Structured logging, Kibana dashboards, Grafana metrics \\
\bottomrule
    \end{tabularx}
    \captionof{table}{Sprint planning overview}
    \label{tab:sprint-overview}
\end{center}

\section{Sprint Goal}

Establish the project foundation, configure the development environment, and implement the first pipeline operator: RabbitMQ message consumption.

\section{Development Environment}

\subsection{Software Environment}
This section lists all the software and tools used throughout the project implementation.

\begin{itemize}
    \item[-] \textbf{IntelliJ IDEA}~\cite{Intel}, is an integrated development environment designed for software development based on Java technology.
    \item[-] \textbf{Minikube}~\cite{minikube}, is an open-source tool that facilitates setting up a local Kubernetes cluster on your machine. It allows developers to test, develop, and deploy Kubernetes applications without requiring a full Kubernetes cluster infrastructure.
    \item[-] \textbf{Docker-cli}, is the command-line interface that allows users to interact with Docker. It enables users to perform a variety of tasks related to managing containers, images, networks, and Docker volumes.
    \item[-] \textbf{Kubectl}, is a Kubernetes command-line tool. It allows executing commands in Kubernetes clusters. It is used to deploy applications, inspect and manage cluster resources, and view logs.
    \item[-] \textbf{Helm}~\cite{helm}, is an open-source package management tool for Kubernetes. It simplifies and automates the deployment of applications and services on Kubernetes clusters using configuration files called charts.
\end{itemize}

\subsection{Technology Stack}

\begin{center}
    \justifying
    \begin{tabularx}{\textwidth}{|l | l | X|}
\toprule
\textbf{Category} & \textbf{Technology} & \textbf{Version / Details} \\
\midrule
Language & Java & 17 (LTS) \\
\hline
Streaming Framework & Apache Flink~\cite{FLINK} & 1.20.2 \\
\hline
Build & Apache Maven & Shade plugin (fat JAR) \\
\hline
Message Broker & RabbitMQ~\cite{RBMQ} & Flink Connector 3.0.1-1.17 \\
\hline
State Store & Valkey~\cite{valkey} & REST API (Redis-compatible) \\
\hline
Notification API & Erable & HTTP/REST via Spring WebClient \\
\hline
AI Inference & ONNX Runtime~\cite{onnxruntime} & 1.17.0 \\
\hline
Serialization & Jackson & 2.15.2 \\
\hline
Validation & Hibernate Validator & 6.2.5.Final (Bean Validation 2.0) \\
\hline
Logging & Logback + Logstash encoder & 1.2.13 / 7.2 \\
\hline
Containerization & Docker~\cite{docker} & Based on Flink 1.20.1 image \\
\hline
Orchestration & Kubernetes~\cite{K8S} / Docker Compose & High availability with ZooKeeper \\
\hline
Authentication & OKAPI v2 & OAuth2 token \\
\bottomrule
    \end{tabularx}
    \captionof{table}{LFU job technology stack}
    \label{tab:stack}
\end{center}

\section{Project Structure}

The \texttt{pushavoo-flink} project is organized into Maven modules:

\begin{itemize}
    \item \texttt{pushavoo-flink-core} --- Shared library (utilities, metrics, RabbitMQ connectors, Erable client, Valkey access)
    \item \texttt{pushavoo-flink-jobs/job-fair-use} --- The LFU job (subject of this report)
    \item \texttt{pushavoo-flink-jobs/job-fiber-eligibility} --- Fiber eligibility job
    \item \texttt{pushavoo-flink-jobs/job-contracts-portfolio-changes} --- Portfolio changes job
    \item \texttt{pushavoo-flink-jobs/job-optimal-renewal-date} --- Optimal renewal date job
    \item \texttt{pushavoo-flink-test} --- Shared test library
\end{itemize}

The \texttt{job-fair-use} module directory structure is as follows:

\begin{verbatim}
job-fair-use/
├── pom.xml
└── src/main/java/.../lfu/
    ├── LFUJob.java                          (Entry point)
    ├── LFUJobConstants.java
    ├── model/
    │   ├── EnrichedFairUseMessage.java
    │   ├── FairUseMessage.java
    │   ├── FairUse.java
    │   ├── LFUErableEntity.java
    │   └── MessageRules.java
    ├── function/
    │   ├── LFUCheckValidMessageFlatMapFunction.java
    │   ├── LFUGetSubscriptionProcessFunction.java
    │   ├── LFUErableCallProcessFunction.java
    │   ├── LFUUpdateSubscriptionSink.java
    │   └── ai/
    │       ├── LFUPredictionProcessFunction.java
    │       └── PredictionFeatureExtractor.java
    └── utils/
\end{verbatim}

\section{Data Model}

\subsection{Input Message: FairUseMessage}

\begin{center}
    \justifying
    \begin{tabularx}{\textwidth}{|l | l | l | X|}
\toprule
\textbf{Field} & \textbf{Type} & \textbf{Constraint} & \textbf{Description} \\
\midrule
\texttt{msisdn} & String & @NotNull, @ValidFrenchMSISDN & Customer phone number (French format) \\
\hline
\texttt{fairuse} & FairUse & @NotNull, @Valid & Fair-use consumption data \\
\bottomrule
    \end{tabularx}
    \captionof{table}{FairUseMessage structure}
    \label{tab:fairusemessage}
\end{center}

\begin{center}
    \justifying
    \begin{tabularx}{\textwidth}{|l | l | p{4cm} | X|}
\toprule
\textbf{Field} & \textbf{Type} & \textbf{Values} & \textbf{Description} \\
\midrule
\texttt{customerType} & String & GP, ENT, MVNO & Customer type (General Public, Enterprise, MVNO) \\
\hline
\texttt{reachedThresholdPercent} & Integer & [0, 99999] & Percentage of quota reached \\
\hline
\texttt{eventType} & String & fairuse, blocage, redirect, alerting, leveeFU & Consumption event type \\
\hline
\texttt{eventDate} & Date & ISO 8601 & Event date and time \\
\bottomrule
    \end{tabularx}
    \captionof{table}{FairUse object structure}
    \label{tab:fairuse}
\end{center}

\subsection{Enriched Message: EnrichedFairUseMessage}

Throughout the pipeline, the message is progressively enriched:

\begin{center}
    \justifying
    \begin{tabularx}{\textwidth}{|l | l | X|}
\toprule
\textbf{Field} & \textbf{Type} & \textbf{Added by} \\
\midrule
\texttt{uuid} & String & Operator 2 --- Validation (Sprint 2) \\
\hline
\texttt{fairUseMessage} & FairUseMessage & Operator 2 --- Validation (Sprint 2) \\
\hline
\texttt{lfuMainRecord} & LfuMainRecord & Operator 3 --- Valkey subscription (Sprint 2) \\
\hline
\texttt{predictionScore} & double & Operator 4 --- AI prediction (Sprint 4) \\
\hline
\texttt{anomalyDetected} & boolean & Operator 4 --- AI prediction (Sprint 4) \\
\bottomrule
    \end{tabularx}
    \captionof{table}{Enriched message structure}
    \label{tab:enriched}
\end{center}

\section{RabbitMQ Source Operator}

RabbitMQ~\cite{RBMQ} plays a central role as the primary data source for our Flink job. The source consumes JSON messages from the queue configured via the \texttt{RMQ\_LFU\_QUEUE} environment variable. It uses a custom RabbitMQ connector (\texttt{CustomRmqSource}) that optionally supports \textit{correlation ID} for message tracking.

It is important to note that the connection to RabbitMQ is made via the secure AMQPS protocol, guaranteeing a secure connection and reinforcing the reliability and confidentiality of data exchanges.

\begin{lstlisting}[language=Java, caption={RabbitMQ source creation}]
CustomRmqSource<String> source = RmqUtils.createCustomRmqSource(
    RMQ_LFU_QUEUE, USE_CORRELATION_ID);
DataStream<String> stream = env.addSource(source)
    .name("Rabbit MQ").uid("rabbit-mq").rebalance();
\end{lstlisting}

\section*{Conclusion}

At the end of Sprint~1, the project foundation was established with a fully configured Maven multi-module structure, development environment, and technology stack. The first pipeline operator---the RabbitMQ source---was implemented, enabling the job to consume raw JSON messages from the queue via the secure AMQPS protocol. Unit tests confirmed correct message consumption and stream creation.


